% answers_short.tex —— 2009 年·普通高等学校招生考试·山东卷·理科 · 简短答案速查
% 由 tools/make_short_answers.py 自 answers.tex 的 \gksolhead/\gkanswer 提取，
% 并与 2009/2009_数学真题_山东理_答案要点.md 逐题比对后生成，勿手改（改源 MD 后重跑）。

\begin{gknotice}
\textbf{使用说明}：本节只给最终答案便于快速核对，完整推导见解析版。答案为编写者独立推导，\textbf{非官方标准答案}。
\end{gknotice}

\gkpart{选择题}
\begin{center}\begin{tabular}{ccccccccccccc}
\toprule
题号 & 1 & 2 & 3 & 4 & 5 & 6 & 7 & 8 & 9 & 10 & 11 & 12 \\
\midrule
答案 & D & C & B & C & B & A & B & A & D & C & A & A \\
\bottomrule
\end{tabular}\end{center}

\gkpart{填空题}
\begin{itemize}[leftmargin=2.6em]
  \item[\textbf{13.}] $(-1,1)$。
  \item[\textbf{14.}] $(1,+\infty)$。
  \item[\textbf{15.}] $30$。
  \item[\textbf{16.}] $-8$。
\end{itemize}

\gkpart{解答题}
\begin{itemize}[leftmargin=2.6em]
  \item[\textbf{17.}] （1）最大值 $\dfrac{1+\sqrt3}{2}$，最小正周期 $\pi$；（2）$\dfrac13$。
  \item[\textbf{18.}] （1）证明见解析；（2）$\dfrac{\sqrt7}{7}$。
  \item[\textbf{19.}] （1）$q_2=0.8$；（2）$E\xi=3.63$；（3）都在 $B$ 处投的得分超过 $3$ 分的概率更大。
  \item[\textbf{20.}] （1）$r=-1$；（2）证明见解析。
  \item[\textbf{21.}] （1）$y=\dfrac{4}{x^2}+\dfrac{9}{400-x^2}$（$0<x<20$）；（2）存在，该点到城 $A$ 的距离为 $4\sqrt{10}\ \mathrm{km}$。
  \item[\textbf{22.}] （1）$\dfrac{x^2}{8}+\dfrac{y^2}{4}=1$；（2）存在，圆 $x^2+y^2=\dfrac83$，$|AB|\in\left[\dfrac{4\sqrt6}{3},2\sqrt3\right]$。
\end{itemize}
